% !TeX program = xelatex
% Standalone: no external assets or shell escape. Data cutoff: 2026-10-17.
% Repository reference: myx / 5d349e82.
\documentclass[UTF8,aspectratio=169,10pt,fontset=fandol]{ctexbeamer}
\usepackage{booktabs}
\usepackage{listings}
\setmonofont{Latin Modern Mono}

\usetheme{default}
\usefonttheme{professionalfonts}
\setbeamertemplate{navigation symbols}{}
\setbeamertemplate{headline}{}
\setbeamertemplate{itemize item}{\raisebox{0.15ex}{\scriptsize$\bullet$}}
\setbeamercolor{normal text}{fg=black,bg=white}
\setbeamercolor{structure}{fg=black}
\setbeamercolor{frametitle}{fg=black,bg=white}
\setbeamercolor{footline}{fg=black,bg=white}
\setbeamerfont{frametitle}{size=\large,series=\bfseries}
\setbeamersize{text margin left=8mm,text margin right=8mm}
\setlength{\parskip}{1pt}

\newcommand{\code}[1]{\texttt{\detokenize{#1}}}
\newcommand{\PerfNote}{内部单调计时包含首次调用；无计时前 workload 预跑。}
\newcommand{\DiagNote}{诊断数据，不作正式性能结果。}
\newcommand{\PageSource}{}
\newcommand{\PageMode}{}
\newcommand{\PageMeta}[2]{\gdef\PageSource{#1}\gdef\PageMode{#2}}
% Fixed footer area: source text is not pushed below a frame's body.
\setbeamertemplate{footline}{%
  \begin{beamercolorbox}[wd=\paperwidth,leftskip=8mm,rightskip=8mm]{footline}
    \setlength{\parskip}{0pt}
    \vskip1pt
    {\fontsize{6.5}{7.5}\selectfont
      来源：\PageSource\hfill\insertframenumber/\inserttotalframenumber\par
      WSL2 Ubuntu24.04.2/x86\_64；QuickJS 2026-06-04；V8 15.6.21（shell: d8）。\par
      V8 \code{--max-opt=0 --no-lazy}。\PageMode\par}
    \vskip2pt
  \end{beamercolorbox}}

\lstdefinestyle{codeblock}{
  basicstyle=\ttfamily\fontsize{9.4}{10.6}\selectfont,
  keywordstyle=\color{black},commentstyle=\color{black},stringstyle=\color{black},
  columns=fullflexible,keepspaces=true,showstringspaces=false,
  breaklines=true,breakatwhitespace=true,
  postbreak=\mbox{\scriptsize$\hookrightarrow$\space},
  aboveskip=2pt,belowskip=2pt,frame=none,tabsize=2}
\lstset{style=codeblock}
\title{QuickJS与V8：字节码和属性访问分析}
\author{马云霄}
\date{} % The meeting date is unknown; do not invent one.

\begin{document}

% -------------------- Main 1/6 --------------------
\PageMeta{S1：numeric\_loop 源码、真实 dump、summary.csv。}{\PerfNote}
\begin{frame}[fragile,t]{QuickJS与V8：字节码和属性访问分析}
\small 马云霄｜截至2026-10-17：本周做小程序字节码对照与原 Richards 属性分析。

\medskip
\textbf{先看数值求和：同一轮循环怎样执行？}
\begin{columns}[T,onlytextwidth]
  \begin{column}{0.33\textwidth}
    \textbf{相同 JavaScript}
\begin{lstlisting}
function wbNumericLoop(n) {
  var sum = 0;
  for (var i=0; i<n; i++) {
    sum += i;
  }
  return sum;
}
\end{lstlisting}
    \small 输入 \code{n=4096}；\par
    返回 8386560，未溢出 int32。
  \end{column}
  \begin{column}{0.24\textwidth}
    \textbf{QuickJS：循环体摘录}
\begin{lstlisting}
get_loc1 1: i
add_loc 0: sum
inc_loc 1: i
goto8 4
\end{lstlisting}
    \small \code{add_loc} 更新局部\par
    \code{sum} 并弹栈；\par
    \code{inc_loc} 更新 \code{i}。
  \end{column}
  \begin{column}{0.40\textwidth}
    \textbf{Ignition：同一段摘录}
\begin{lstlisting}
Ldar r1
Add r0, EmbeddedFeedback[None]
Mov r0, r2
Star0
Ldar r1
Inc EmbeddedFeedback[None]
Star1
JumpLoop [21], [0], FBV[0]
\end{lstlisting}
    \small \code{r0=sum}，\code{r1=i}；\code{Star0} 写回。\par
    \code{Mov} 复制旧 sum；虚拟槽非 CPU 寄存器。
  \end{column}
\end{columns}
\smallskip
\textbf{Median（$\mu$s/完整调用）：}
QJS 43.782；V8 54.839；V8/QJS 1.253。
\par\small 每完整目标调用含4096轮；每引擎30个新进程，点估计。
\raggedright
\small \textbf{本例 QJS 更快}；指令条数/字节数不等于成本，也不能据此断言栈式更优。
\end{frame}

% -------------------- Main 2/6 --------------------
\PageMeta{S2：四组 cases、summary.csv、raw/timing.csv。}{\PerfNote}
\begin{frame}[fragile,t]{四组调用方式：一层原型够不够解释差距？}
共同计算：\code{i=0..4095}、\code{bias=3}，累加 \code{i+bias}，结果8398848。
\smallskip
\centering
\setlength{\tabcolsep}{5pt}
\renewcommand{\arraystretch}{1.22}
\begin{tabular}{@{}llrrr@{}}
  \toprule
  组别 & 循环中的主要表达式 & QJS / $\mu$s & V8 / $\mu$s & V8/QJS \\
  \midrule
  直接计算 & \code{sum += i + obj.bias} & 67.132 & 60.388 & 0.900 \\
  普通调用 & \code{sum += stepPlain(obj,i)} & 159.969 & 116.395 & 0.728 \\
  own 方法 & \code{sum += obj.step(i)} & 217.156 & 120.905 & \textbf{0.557} \\
  prototype 方法 & \code{sum += obj.step(i)} & 225.086 & 125.348 & \textbf{0.557} \\
  \bottomrule
\end{tabular}
\par\smallskip
\raggedright
\small 中位时间单位：$\mu$s/完整目标调用，含4096轮内循环及摊薄 driver；不是单次操作成本。\par
每组每引擎30个新进程，共240个正式样本；对象建立在计时外，方法读取留在循环内。
\medskip
\begin{columns}[T,onlytextwidth]
  \begin{column}{0.49\textwidth}
    \textbf{普通函数与方法的计算体}
\begin{lstlisting}
function stepPlain(obj, x) {
  return x + obj.bias;
}
function stepMethod(x) {
  return x + this.bias;
}
\end{lstlisting}
    \small own/prototype 复用相同循环和方法；\par
    只在计时外改变方法所属对象。
  \end{column}
  \begin{column}{0.48\textwidth}
    \textbf{方法变为 own 后，跨引擎差距仍在。}\par
    一层原型不足以解释这组差距。\par\smallskip
    参数、receiver、对象布局/反馈也有差异；\par
    组间相减不能叫纯调用或纯查找成本。\par\smallskip
    \small 本页直接计算含 \code{obj.bias} 与两次加法；\par
    \textbf{不是第1页的单纯求和。}
  \end{column}
\end{columns}
\end{frame}

% -------------------- Main 3/6 --------------------
\PageMeta{S3：原 richards.js；site\_samples.csv、sampling\_summary.json。}{\DiagNote}
\begin{frame}[fragile,t]{原始 Richards：查找区域样本落在哪里？}
Richards 是模拟任务调度的 JavaScript 基准；原算法、方法读取和断言均未改。
\medskip
\begin{columns}[T,onlytextwidth]
  \begin{column}{0.54\textwidth}
    \textbf{Scheduler.schedule 摘录}\par
    \small 此处 \code{this} 是 Scheduler。
\begin{lstlisting}
this.currentTcb = this.currentTcb.run();
\end{lstlisting}
    \smallskip
    \normalsize\textbf{TaskControlBlock.run 摘录}\par
    \small 进入此函数后，\code{this} 是当前 TCB。
\begin{lstlisting}
return this.task.run(packet);
\end{lstlisting}
    \smallskip
    \normalsize \code{currentTcb} 指向 TCB；\par
    TCB 的 \code{task} 指向具体任务对象。\par
    取得对象、读方法、调用函数不是同一步。
  \end{column}
  \begin{column}{0.43\textwidth}
    \textbf{按属性聚合的查找区域样本}
    \smallskip
    \centering
    \begin{tabular}{@{}lrr@{}}
      \toprule
      属性 & 样本 & 占1229 \\
      \midrule
      \code{currentTcb} & 280 & 22.8\% \\
      \code{state} & 204 & 16.6\% \\
      两处 \code{run} & 79 & 6.4\% \\
      \bottomrule
    \end{tabular}
    \par\smallskip
    \raggedright
    \small 3个插桩诊断进程，共5393个内部样本；\par
    其中1229个查找区域样本已映射。\par\smallskip
    同名属性可含多个读/写位置。\par
    \textbf{分母不是整个程序耗时。}
  \end{column}
\end{columns}
\medskip
\textbf{普通状态/调度字段值得继续看；尚未找到总体瓶颈。}
\par\smallskip
\small 方法说明：189个活跃 opcode 位置；N1/N2计数倍增核对，计数与采样分开。\par
另有39次 helper 探查无具体 site；不把访问次数当作耗时排名。
\end{frame}

% -------------------- Main 4/6 --------------------
\PageMeta{S3：richards.js:309--311；两引擎 isHeldOrSuspended 最终 dump。}{\DiagNote}
\begin{frame}[fragile,t]{选定 this.state：按源码语义分组读字节码}
\small 原函数：\code{TaskControlBlock.prototype.isHeldOrSuspended}。
\begin{lstlisting}
return (this.state & STATE_HELD) != 0 || (this.state == STATE_SUSPENDED);
\end{lstlisting}
\small 调用者 \code{schedule} 的 \code{this} 是 Scheduler；进入本函数后是当前 TCB。
\par ①取得 this/读 state；②按位与、比较、短路；③必要时再读。
\begin{columns}[T,onlytextwidth]
  \begin{column}{0.48\textwidth}
    \normalsize\textbf{QuickJS}
    \par\small ①offset 0--3；②8--15；③18--19
\begin{lstlisting}
 0 push_this
 1 put_loc0 0: this
 2 get_loc0 0: this
 3 get_field state
 8 get_var 0: STATE_HELD
11 and
13 neq
15 if_true8 28
18 get_loc0 0: this
19 get_field state
\end{lstlisting}
  \end{column}
  \begin{column}{0.49\textwidth}
    \normalsize\textbf{V8 Ignition}
    \par\small ①offset 0--4；②5--17；③19
\begin{lstlisting}
 0 GetNamedProperty <this>, [0:"state"], FBV[0]
 4 Star0
 5 LdaGlobal [1:"STATE_HELD"], FBV[2]
 8 BitwiseAnd r0, EmbeddedFeedback[None]
13 TestEqual r0, EmbeddedFeedback[None]
16 LogicalNot
17 JumpIfTrue [13] (0x1eef0100166e @ 30)
19 GetNamedProperty <this>, [0:"state"], FBV[4]
\end{lstlisting}
  \end{column}
\end{columns}
\small \textbf{非连续摘录}，完整见附录；折行不是新指令，\code{r0} 是虚拟槽，offset 不是 native PC。
\end{frame}

% -------------------- Main 5/6 --------------------
\PageMeta{S3：counts\_N1.sites.csv、IC transitions.csv、实际 handler 反汇编。}{\DiagNote}
\begin{frame}[t]{这一处实际知道了什么？}
\begin{columns}[T,onlytextwidth]
  \begin{column}{0.48\textwidth}
    \textbf{QJS：首次 state，offset 3}\par\medskip
    每 Run \textbf{10671次}属性操作：
    \begin{itemize}
      \setlength{\itemsep}{4pt}
      \item 全部 own 普通数据属性命中；\par
            每次一次自身探查，结果为整数。
      \item 原型命中0，通用 fallback 0；\par
            \textbf{不等于读取没有成本。}
      \item 仍有 shape/hash 查找、描述符检查、\par
            取值及必要的 receiver 引用处理。
    \end{itemize}
    \small 整数 state 本身不做对象引用计数；\par
    不能把 receiver 与字段值混为一谈。
  \end{column}
  \begin{column}{0.49\textwidth}
    \textbf{V8：同进程字节码与 IC 日志}\par\smallskip
    \centering
    \begin{tabular}{@{}ccc@{}}
      \toprule
      offset & 反馈槽 & 转换码 \\
      \midrule
      0 & \code{FBV[0]} & $0\to1$ \\
      19 & \code{FBV[4]} & $0\to1$ \\
      \bottomrule
    \end{tabular}
    \par\smallskip
    \raggedright
    相同 receiver map：\code{0x189b01023ef1}。
    \begin{itemize}
      \setlength{\itemsep}{4pt}
      \item 转换日志不是逐次命中记录；\par
            IC命中率仍 UNKNOWN。
      \item handler 有 map 检查/字段读取路径，\par
            不证明此 site 每次走该路径。
    \end{itemize}
    \small 精确源码/构建与制品的对应仍有缺口。
  \end{column}
\end{columns}
\medskip
\textbf{待验证解释：}QJS 重复自身查找；V8 可能复用稳定字段反馈。\par
这是可继续检验的路径差异，\textbf{不是已经证明的性能原因}，也未量出各步骤时间。
\end{frame}

% -------------------- Main 6/6 --------------------
\PageMeta{S1/S2 正式汇总；S3 perturbation.json、site\_samples.csv、N1计数。}{\DiagNote}
\begin{frame}[t]{本周能说什么，还有什么不懂？}
\textbf{目前能说}
\begin{itemize}
  \setlength{\itemsep}{2pt}
  \item 操作不同，快慢方向可能不同；不能给小程序拼一个总体引擎排名。
  \item 方法案例：多一层原型不足以解释跨引擎差距。
  \item 原 Richards 的普通 own 字段值得深入；已选出可具体分析的 state 位置。
\end{itemize}
\smallskip
\textbf{两个直接限制（均为诊断，不是正式性能结果）}
\begin{itemize}
  \setlength{\itemsep}{3pt}
  \item 无信号小对照，N512：冻结版median1718ms，诊断开启版2223ms，\par
        高约\textbf{29\%}。不同批次点估计，不是精确的插桩成本。
  \item 两次 state 读取每 Run 为10671/10115次，相关查找区域样本却为\textbf{85/9}。\par
        请求1ms，实际约10ms；采样相位与插桩会影响分布，不能当精确耗时排名。
\end{itemize}
\medskip
\textbf{想与老师讨论的问题}\par
下一步先确认 V8 该位置的实际快路径，\par
还是先降低采样扰动、复核位置分布？
\end{frame}

\appendix
% -------------------- Appendix A --------------------
\PageMeta{S3：QuickJS 完整目标函数 dump 与 decoded.json；见 notes.md。}{\DiagNote}
\begin{frame}[fragile,t]{附录A：QuickJS 完整最终指令序列}
\small 16条/29B；局部0=this，栈深2。补齐offset，保留全部指令/操作数，去掉源码注释。
\begin{lstlisting}
 0 push_this
 1 put_loc0 0: this
 2 get_loc0 0: this
 3 get_field state
 8 get_var 0: STATE_HELD
11 and
12 push_0 0
13 neq
14 dup
15 if_true8 28
17 drop
18 get_loc0 0: this
19 get_field state
24 get_var 1: STATE_SUSPENDED
27 eq
28 return
\end{lstlisting}
\end{frame}

% -------------------- Appendix B --------------------
\PageMeta{S3：同进程 V8 完整目标函数 dump；IC地址映射见 notes.md。}{\DiagNote}
\begin{frame}[fragile,t]{附录B：Ignition 完整最终指令序列}
\small 14条/31B；虚拟寄存器1、frame 8B。去掉地址/编码/调试列，指令与操作数完整。
\begin{lstlisting}
 0 GetNamedProperty <this>, [0:"state"], FBV[0]
 4 Star0
 5 LdaGlobal [1:"STATE_HELD"], FBV[2]
 8 BitwiseAnd r0, EmbeddedFeedback[None]
11 Star0
12 LdaZero
13 TestEqual r0, EmbeddedFeedback[None]
16 LogicalNot
17 JumpIfTrue [13] (0x1eef0100166e @ 30)
19 GetNamedProperty <this>, [0:"state"], FBV[4]
23 Star0
24 LdaGlobal [2:"STATE_SUSPENDED"], FBV[6]
27 TestEqual r0, EmbeddedFeedback[None]
30 Return
\end{lstlisting}
\small QJS 16条/29B，V8 14条/31B；条数、编码长度不是执行时间。
\end{frame}

\end{document}
